% Paper Summary: WebGPT: Browser-Assisted Question-Answering with Human Feedback (Nakano et al., 2021)

\begin{papersummary}{2021}{WebGPT: Browser-Assisted Question-Answering with Human Feedback}{Reiichiro Nakano, Jacob Hilton, Suchir Balaji, Jeff Wu, Long Ouyang, Christina Kim, Christopher Hesse, Shantanu Jain, Vineet Kosaraju, William Saunders, Xu Jiang, Karl Cobbe, Tyna Eloundou, Gretchen Krueger, Kevin Button, Matthew Knight, Benjamin Chess, John Schulman}{WebGPT fine-tuned GPT-3 to answer long-form questions by issuing search queries and navigating a text-based browser, trained with human feedback to prefer well-supported, properly cited answers over unsupported ones.}

\summaryheading{Key Ideas}
A language model's parametric knowledge is fixed at training time and unverifiable at inference time; WebGPT instead gives the model a text-based browsing environment---search, click links, scroll, quote passages---and trains it to collect supporting references while composing an answer. Training proceeds in the now-familiar three stages: behavior cloning from human demonstrations of browsing and answering, a reward model trained on human comparisons between candidate answers, and policy optimization against that reward model. Human evaluators preferred WebGPT's answers, complete with citations to the pages it visited, over unsupported model answers and, on a meaningful fraction of comparisons, over human-written ones.

\summaryheading{Follow-on Works}
The pattern of grounding a language model's actions in an external environment and training the interaction with human feedback recurs throughout this appendix: \hyperref[paper:schick-2023]{Toolformer (Schick et al., 2023)} generalized from a single browsing tool to arbitrary APIs the model learns to invoke on its own, and \hyperref[paper:jin-2025-search-r1]{Search-R1 (Jin et al., 2025)} replaced WebGPT's imitation-and-reward-model pipeline with reinforcement learning against verifiable outcome rewards, training search behavior directly rather than through a learned proxy for human preference. The RLHF machinery WebGPT applied to browsing was, in the same period, applied to general instruction-following in \hyperref[paper:ouyang-2022]{InstructGPT (Ouyang et al., 2022)}, sharing several authors and the same three-stage training recipe.

\summaryheading{Lasting Contributions}
WebGPT is the origin point of tool-augmented, human-feedback-trained language agents: the idea that a model's actions in an external environment, not just its text output, can be shaped by the same demonstration-and-preference pipeline used for chat alignment. Citation-grounded, browsing-capable question answering is now a standard feature of deployed assistants, and the retrieval-augmented, tool-use lineage that runs through this appendix to Search-R1 traces back to this paper's specific combination of a real (if simplified) web environment with human preference training, four years before that combination became common practice.

\end{papersummary}
